\documentclass[11pt]{article}
\usepackage[margin=1in]{geometry}
\usepackage{hyperref}

\title{H3 Result Memo (2026)}
\author{C.M. Baird et al.}
\date{March 2026}

\begin{document}
\maketitle

\section*{Purpose}
This memo summarizes the current H3 executable joint posterior / model-comparison package in publication-facing language, with explicit interpretation boundaries.

\section*{Hypothesis}
H3 asks whether one shared parameter vector remains viable across:
\begin{itemize}
    \item H1 simulation-primary branch-selection constraints,
    \item H2 interferometric visibility sensitivity outputs,
    \item fifth-force bound tables.
\end{itemize}

\section*{Current executable status}
The current H3 package is no longer only a validity scaffold. It now includes:
\begin{itemize}
    \item a prototype joint posterior sampler,
    \item a null-model comparison score,
    \item leave-one-channel-out diagnostics.
\end{itemize}

\section*{Current outputs}
Using the present implemented channel set, the prototype returns:
\begin{itemize}
    \item posterior mean $\eta \approx 1.66\times 10^{-2}$,
    \item posterior mean $\gamma \approx 2.43\times 10^{5}$,
    \item posterior mean $\lambda_{\phi} \approx 2.07\times 10^{-12}$,
    \item posterior mean $m_{\phi} \approx 9.61\times 10^{3}$,
    \item log-evidence-like advantage over null of approximately $15.38$ in the current prototype scoring convention.
\end{itemize}

\section*{Interpretation}
This result means the \emph{currently implemented} H1+H2+fifth-force package admits a nontrivial jointly weighted region under the present prototype assumptions. It does \textbf{not} yet mean that nature prefers the theory, nor that the full collider/cosmology channel set has been incorporated at publication-grade level.

\section*{Boundaries}
The current H3 result should be stated only as:
\begin{itemize}
    \item an executable multi-channel inference prototype,
    \item a positive internal consistency result under the current implemented channels,
    \item a basis for next-stage joint-fit strengthening.
\end{itemize}

It should \textbf{not} yet be stated as:
\begin{itemize}
    \item decisive external evidence for MQGT-SCF,
    \item a finished global posterior over all intended channels,
    \item a replacement for independent external replication.
\end{itemize}

\section*{Next scientific step}
The immediate next move is to extend the H3 posterior package to include the missing collider/cosmology tables and maintain the same bounded claim style.

\section*{Bottom line}
H3 has advanced from a narrative consistency claim to a real executable inference layer. That is scientifically meaningful progress, but it is still an internal inference milestone rather than a final empirical verdict.

\end{document}
